\section{合卷：霸权为什么能维持，又为什么不能变成帝国}

齐、晋、楚的霸权都依赖联盟、礼仪承认和军事威慑。它们能在某个阶段压住冲突，却难以把其他诸侯变成稳定的行政单位。会盟是一种低成本秩序：霸主不必在每个国家驻官、征税或养军，只要能在关键时刻召集盟友、裁决争端、保护小国，秩序就能暂时运转；代价是秩序必须不断重新谈判。

\subsection{“霸”不是一把可以传下去的椅子}

后世将齐桓、晋文、楚庄并列，很容易以为他们先后坐上同一个位置。春秋材料所见的却不是一项由周廷授予、可以按期交接的职位。每一个能居于盟坛中央的国家，都先要在具体危机中证明：它能否使援军越过道路，能否让城邑愿意供给，能否让受威胁的小国相信，到场或出兵比拒绝更安全。名分使这种证明容易被旁人读懂，不能替它完成。

齐的路径最能说明这一点。北杏之会中宋的拒绝，\eventref{SA-0681-BEIXING}{春秋·北杏之会}已经显示会盟并非命令发出便自然生效；齐此后援燕、迁邢，尤其是\eventref{SA-0659-QI-XING}{春秋·齐迁邢于夷仪}，才把一个强国能够为边患中的小国出兵、安置和续接关系的能力放到列国眼前。到\eventref{SA-0656-SHAOLING}{春秋·召陵之盟}，齐军及其盟友能够抵达楚北、以盟收束冲突；\eventref{SA-0651-KUIQIU}{春秋·葵丘之会}又表明这张网络尚能被再次召集。它们不是齐已经统治蔡、陈、江、黄或楚的证据，而是齐在不同地点把保护、压力与周礼语言暂时接合的几次成功。

援助并不只是强国施恩。迁徙者要离开旧田与旧城，参与会盟的小国要提供使者、通道或军粮，齐也要为远征承担本国的征发与风险。受援者可从齐的组织中取得喘息，日后也会被期待响应齐的召集；这种互相期待使关系有了延续性，却没有取消各自重新选择的权利。齐桓死后的继嗣冲突之所以迅速削弱其号令，正因为盟会所倚赖的中心首先无法在本国平稳地交接。

\subsection{霸主靠什么取得位置}

\eventref{SA-0685-QIHUAN}{春秋·齐桓公即位}后的齐，靠海滨资源、内部整合和“尊王攘夷”把本国利益说成公共秩序；\eventref{SA-0632-CHENGPU}{春秋·城濮之战}后的晋，靠重耳流亡时结成的关系、北方腹地和军事联盟成为中原枢纽；\eventref{SA-0597-BI}{春秋·邲之战}后的楚，则以江汉腹地和北上能力打破“中原国家自然领先”的想象。

三种霸权的共同点是：它们都需要其他国家的承认。霸主若只会出兵，盟友会担心它变成新的吞并者；若只会讲礼，联盟又会在危机中失去保护。因此霸权总在军事实力与礼仪克制之间摆动。

\subsection{胜利必须再被安放进秩序}

晋文公的经历说明，一场会战本身也不能产生稳定的号令。王子带之乱中，晋助周襄王复归王都，\eventref{SA-0635-XIANG-WANG}{春秋·襄王复归王都}使晋的兵力能够以援王的语言被理解；城濮取胜后，晋又在\eventref{SA-0632-JIANTU}{春秋·践土之会}与周王名分相接。周王室从中得到延续自身的外援，晋则把已经取得的军事优势转换成了召集诸侯的资格。双方是互相借重，不是周廷把中原的军队、粮食和城邑忽然交给晋直接管理。

楚在邲之战后的声望同样来自这种转换是否可被维持。前598年楚围郑、晋援未能阻止郑暂时转向，\eventref{SA-0598-CHU-ZHENG-SIEGE}{春秋·楚围郑}先动摇的是晋方保护承诺的信用；次年的\eventref{SA-0597-BI}{春秋·邲之战}才把这种压力集中为会战。楚军获胜使郑、陈、蔡等国必须重新估量向谁求援，却没有使楚拥有一套能够常设征税、任命中原诸侯的官府。江汉腹地给楚提供了再度北进的条件，离开腹地后的粮道、属国和将领合作又始终限定它能把胜利延长多远。

因此，霸权的核心不宜被缩成“哪国最强”。它是一种把本国资源、外部承认和危机中的行动接到一起的能力。齐把北方救援与南向会盟相接，晋把援王、会战与践土礼仪相接，楚则把江汉资源、围郑与北方属国相接；三者在地理、组织与对手上都不同。它们共享的限度是，每一次接合都要有人愿意出兵、开路、供粮或承认，不能由一次胜利永久替代。

\subsection{为什么会盟不能变成常设政府}

齐桓在葵丘、晋文在城濮之后都能调动诸侯，却不能规定每一国如何收税、如何继承、如何任用卿大夫。宋能主持\eventref{SA-0546-MIBING}{春秋·弭兵之会}，正说明大国也需要一个中介；但弭兵没有常设执行机构，齐和秦也没有完全接受它。

这种局面让小国有生存空间，也让大国反复陷入成本高昂的战争。对外会盟需要资源，国内卿族也需要分配资源。晋在霸权最强时把军政交给多支卿族，最后却由这些卿族分走了国家本身。

\subsection{程序能降低成本，不能取消成本}

晋在前六世纪中叶反复尝试把临时动员做得较可预期。邢丘之会试图规整朝聘次序，\eventref{SA-0565-XINGQIU}{春秋·邢丘之会}所见的不是一套保存完整的外交章程，而是以到场、先后与礼仪把对郑、蔡等地的压力变成可重复程序的努力。次年伐郑并盟于戏，\eventref{SA-0564-XI-ALLIANCE}{春秋·伐郑与戏盟}又立刻露出程序背后的强制：不愿稳定站入晋方网络的国家，仍会面对军队与征发。会盟不是战争的反面；它是让战争不必每次都从零开始的一种办法。

新君继位时，这种办法尤其脆弱。晋悼公死后，溴梁由大夫会盟，\eventref{SA-0557-MULIANG-ALLIANCE}{春秋·溴梁之盟}提示掌握军政、使节与随从的人已不能只被看作国君的无声代表。它不证明“大夫取代国君”，也不说明所有诸侯已经形成同一种官僚结构；它只说明承诺由谁背负、谁能把承诺带回城邑落实，成了联盟存续的实际问题。齐晋战后在澶渊会盟，\eventref{SA-0553-CHANYUAN}{春秋·澶渊之盟}能够暂停公开冲突，却没有保存一份足以强制未来履约的条约文本。

宋向戌调停的弭兵之会因而不是春秋国际制度的完成式。长期战争使晋、楚及居间诸侯都承受人力、道路与供给的消耗，宋尚能充当中介，才使\eventref{SA-0546-MIBING}{春秋·弭兵之会}成为可能。它暂时承认两大网络各有从属关系，降低了直接决战的频率；但齐、秦未全然纳入，郑、陈等地的处境也并未从此固定。更重要的是，外部战事缓和后，军将职位、封邑与家臣网络没有回到公室，反而有更多时间在各家门中累积。会盟解决了眼前的互相消耗，却不能替晋、楚或鲁解决国内资源归谁调度的问题。

\begin{sourcenote}[盟坛留下了什么，又没有留下什么]
《春秋》最能固定北杏、召陵、践土、邢丘、溴梁、澶渊与弭兵等行动的鲁历年次、参会者和极简结果；它通常不保存盟辞、兵粮数额或违约的处置。那些更完整的使者言辞、礼仪次序和人物判断，多出自《左传》与《国语》的叙事层；《史记》则把它们组织进汉代所见的霸业兴亡线。因而会盟可被当作关系正在被公开安排的证据，不能被误读为一套留有全文、覆盖列国的国际法。历史地理研究能够追问会地、道路和城邑的约束，也不能据此替每一国补出未存的供给簿。
\end{sourcenote}

\subsection{从霸权的得失到战国的难题}

\begin{turningpoint}[制度得失：低成本秩序的边界]
\term{得}：会盟、朝聘与礼仪让缺乏统一官僚体系的诸侯世界仍能协调战争、救援和外交；对小国而言，它提供了向霸主申诉、借力制衡的渠道。

\term{失}：霸权依赖人物、声望和持续资源，无法把联盟转化为可继承的公共制度。每次国君去世、卿族争权、边疆战争或大国换位，旧盟约都需要重新谈判。

\term{后续压力}：战国国家将尝试用县邑、户籍、军功、税收和官吏考核替代会盟的临时协调。它们会得到更强的动员能力，也会让普通家庭更直接地承受国家竞争。
\end{turningpoint}

春秋末年的吴又把这个限度带到更远的空间。前506年，晋方诸侯在召陵侵楚，\eventref{SA-0506-ZHAOLING-INVASION}{春秋·召陵会师侵楚}；同年吴、蔡、唐的军队在\eventref{SA-0506-BOJU}{春秋·柏举之战}攻入郢都。两场行动不可混成同一支“诸侯联军”：前者是晋方重新组织北方关系，后者来自楚东面的吴楚竞争。它们同时压向楚，说明春秋后期的秩序已不是齐、晋、楚轮流主持的一条直线，也说明大国的防务必须同时面对不同方向、不同联盟和不同补给线。

夫差在黄池争取北方盟主次序，\eventref{SA-0482-HUANGCHI}{春秋·黄池之会}却使越得到进攻吴本土的机会；至\eventref{SA-0473-WU-FALL}{春秋·越灭吴}，吴的王室政权退出竞争。吴的兴亡不是对会盟无用的证明：黄池确实能带来声望、联络和可见的政治位置。它说明的是，北方的名分不能替代南方的守备；一国无法把所有资源同时投入远方会盟、对外远征和本土安全。霸权的半径，总小于它声称可以覆盖的世界。

因此，春秋留下的不是统一帝国，而是一套更尖锐的问题：怎样把封国和卿族掌握的资源重新集中？怎样让军役、土地和行政不只依赖贵族关系？这些问题将在战国各国的变法中得到截然不同的回答。
